\documentclass[../DoAn.tex]{subfiles}

\begin{document}

\section{Mục tiêu của luồng ingest}
\label{section:chapter4-ingest-goal}

Chương này trình bày thiết kế luồng ingest, tức luồng tiếp nhận và trích xuất nội
dung từ tài liệu PDF. Trọng tâm của chương là cách hệ thống chuyển một tài liệu
PDF đầu vào thành dữ liệu có cấu trúc để các bước phía sau có thể lập chỉ mục,
truy xuất và dẫn chứng. Đây là phần gắn trực tiếp với tên đề tài, vì chất lượng
trích xuất PDF quyết định dữ liệu mà hệ thống có thể khai thác.

Luồng ingest không chỉ lấy văn bản từ PDF. Hệ thống cần xác định trang, vùng nội
dung, thứ tự đọc, loại block, bảng, ô bảng và các metadata nguồn. Các thông tin
này giúp kết quả trích xuất không bị tách rời khỏi tài liệu gốc, đồng thời cho phép
truy vết lại câu trả lời về đúng trang, đoạn văn hoặc ô bảng.

\section{Tổng quan luồng ingest}
\label{section:ingest-flow-overview}

Luồng ingest được tổ chức thành các bước có thể kiểm tra độc lập. Đầu vào là
một tài liệu PDF; đầu ra là tập block và chunk có metadata, có thể đưa vào chỉ mục
truy xuất. Các bước liên quan đến OCR và trích xuất bảng sử dụng các nền tảng
đã trình bày ở Chương 3, gồm PP-OCR và Table Transformer
\cite{du2020ppocr,smock2022pubtables}.

\begin{figure}[H]
\centering
\small
\begin{tabular}{c}
\fbox{\parbox{0.78\linewidth}{\centering PDF đầu vào}}\\[0.15cm]
\(\downarrow\)\\[-0.05cm]
\fbox{\parbox{0.78\linewidth}{\centering Thăm dò tài liệu: kiểm tra lớp văn bản, ảnh, trang scan, số block}}\\[0.15cm]
\(\downarrow\)\\[-0.05cm]
\fbox{\parbox{0.78\linewidth}{\centering Lập kế hoạch ingest: chọn trích xuất văn bản, bố cục, OCR hoặc định tuyến vùng}}\\[0.15cm]
\(\downarrow\)\\[-0.05cm]
\fbox{\parbox{0.78\linewidth}{\centering Phát hiện vùng nội dung và sắp xếp thứ tự đọc}}\\[0.15cm]
\(\downarrow\)\\[-0.05cm]
\fbox{\parbox{0.78\linewidth}{\centering Định tuyến vùng: văn bản, bảng, ảnh/scan sang bộ xử lý phù hợp}}\\[0.15cm]
\(\downarrow\)\\[-0.05cm]
\fbox{\parbox{0.78\linewidth}{\centering Chuẩn hóa block, tái dựng bảng, sinh chunk và metadata}}\\
\end{tabular}
\caption{Luồng ingest và trích xuất dữ liệu từ PDF}
\label{fig:chapter4-ingest-flow}
\end{figure}

Thiết kế này tách rõ hai trách nhiệm. Probe hoạt động ở mức tài liệu hoặc trang
để quyết định chiến lược xử lý tổng thể. Region routing hoạt động ở mức vùng
nội dung trong trang để chọn backend phù hợp cho từng vùng.

\section{Liên kết yêu cầu với thiết kế ingest}
\label{section:ingest-traceability}

Bảng~\ref{tab:ingest-traceability} tóm tắt cách các yêu cầu liên quan đến trích
xuất PDF được ánh xạ vào thành phần thiết kế. Các mã yêu cầu có thể liên kết với
bảng yêu cầu chức năng ở Chương 2.

\begin{table}[H]
\centering
\small
\renewcommand{\arraystretch}{1.2}
\caption{Liên kết yêu cầu với thành phần thiết kế trong luồng ingest}
\label{tab:ingest-traceability}
\begin{tabularx}{\textwidth}{
    |>{\raggedright\arraybackslash}p{0.16\textwidth}
    |>{\raggedright\arraybackslash}X
    |>{\raggedright\arraybackslash}p{0.30\textwidth}|
}
\hline
\textbf{Mã} & \textbf{Yêu cầu liên quan} & \textbf{Thành phần thiết kế} \\
\hline
FR-01 &
Tiếp nhận tài liệu PDF và phân tích thông tin cơ bản của tài liệu. &
\texttt{probe.py}, luồng điều phối ingest. \\
\hline
FR-02 &
Trích xuất văn bản, bảng và metadata nguồn như trang, vùng nội dung và loại block. &
Phát hiện vùng, sắp xếp thứ tự đọc, chuẩn hóa block, bộ trích xuất văn bản/bảng. \\
\hline
FR-03 &
Định tuyến xử lý theo vùng nội dung như văn bản, bảng, ảnh hoặc vùng scan. &
Định tuyến theo vùng và các bộ xử lý văn bản/bảng/OCR. \\
\hline
FR-04 &
Chia đoạn và lập chỉ mục từ nội dung đã trích xuất. &
Chia đoạn văn bản, chia đoạn có nhận thức bảng, chỉ mục BM25, chỉ mục dense. \\
\hline
FR-07 &
Giữ metadata nguồn để phục vụ dẫn chứng ở các tầng sau. &
Trang, bbox, đường dẫn đề mục, mã bảng, metadata hàng/cột/ô. \\
\hline
\end{tabularx}
\end{table}

\section{Đầu vào và đầu ra của luồng ingest}
\label{section:ingest-input-output}

Bảng~\ref{tab:ingest-input-output} mô tả rõ dữ liệu đi vào và dữ liệu được tạo ra
sau ingest. Phần này giúp tránh nhầm lẫn giữa metadata và dữ liệu trích xuất: mỗi
block/chunk luôn có nội dung chính, còn metadata dùng để mô tả nguồn gốc và
cách xử lý nội dung đó.

\begin{table}[H]
\centering
\small
\renewcommand{\arraystretch}{1.2}
\caption{Dữ liệu vào và đầu ra của luồng ingest}
\label{tab:ingest-input-output}
\begin{tabularx}{\textwidth}{
    |>{\raggedright\arraybackslash}p{0.26\textwidth}
    |>{\raggedright\arraybackslash}X
    |>{\raggedright\arraybackslash}p{0.28\textwidth}|
}
\hline
\textbf{Thành phần} & \textbf{Dữ liệu chính} & \textbf{Mục đích} \\
\hline
PDF đầu vào &
File PDF, số trang, text layer nếu có, ảnh nhúng, hình học trang. &
Nguồn dữ liệu cần trích xuất. \\
\hline
Kết quả probe &
Số trang, số ký tự, số block, số ảnh, tỷ lệ trang có text, tỷ lệ trang có khả năng scan. &
Chọn chiến lược ingest ban đầu. \\
\hline
Region &
Bbox, loại vùng, nội dung thô và backend được định tuyến. &
Chia trang thành các vùng xử lý độc lập. \\
\hline
Block &
Văn bản đã trích xuất, loại block, trang, bbox, heading path và metadata xử lý. &
Đơn vị nội dung sau trích xuất PDF. \\
\hline
Table block &
Caption, bảng Markdown, danh sách cell, hàng, cột, bbox và metadata bảng. &
Giữ cấu trúc bảng để chunking và citation. \\
\hline
Chunk &
Nội dung dùng cho truy xuất, chunk id, trang, loại block và metadata nguồn. &
Đơn vị được đưa vào corpus và retrieval index. \\
\hline
\end{tabularx}
\end{table}

\section{Probe và lập kế hoạch ingest}
\label{section:probe-ingest-plan}

Probe là bước thăm dò nhanh tài liệu trước khi chạy các backend trích xuất. Bước
này không tạo dữ liệu cuối cùng cho truy xuất, mà thu thập đặc trưng để lựa chọn
chiến lược ingest phù hợp. Ví dụ, nếu PDF có text layer rõ ràng, hệ thống ưu tiên
trích xuất văn bản trực tiếp; nếu nhiều trang ít chữ nhưng nhiều ảnh, hệ thống cân
nhắc OCR \cite{du2020ppocr}; nếu trang có bố cục hỗn hợp, hệ thống ưu tiên
region routing.

\begin{table}[H]
\centering
\small
\renewcommand{\arraystretch}{1.2}
\caption{Một số tín hiệu dùng trong bước probe}
\label{tab:probe-signals}
\begin{tabularx}{\textwidth}{
    |>{\raggedright\arraybackslash}p{0.30\textwidth}
    |>{\raggedright\arraybackslash}X|
}
\hline
\textbf{Tín hiệu} & \textbf{Ý nghĩa trong thiết kế} \\
\hline
Tỷ lệ trang có text layer &
Ước lượng khả năng trích xuất văn bản trực tiếp từ PDF. \\
\hline
Tỷ lệ trang không có text &
Gợi ý tài liệu scan hoặc trang ảnh cần OCR. \\
\hline
Số block và mật độ block &
Gợi ý mức độ phức tạp của bố cục trang. \\
\hline
Số ảnh và vùng ảnh lớn &
Gợi ý trang nhiều hình hoặc nội dung cần nhận dạng từ ảnh. \\
\hline
Chất lượng văn bản trung bình &
Giúp phát hiện text layer nhiễu hoặc trích xuất kém. \\
\hline
\end{tabularx}
\end{table}

Sau probe, hệ thống tạo một kế hoạch ingest gồm danh sách backend theo thứ tự ưu
tiên. Nếu backend đầu tiên không tạo được block hợp lệ, cơ chế fallback chuyển
sang backend tiếp theo. Cách này giúp hệ thống không phụ thuộc vào một phương
pháp trích xuất duy nhất.

\section{Định tuyến theo vùng nội dung}
\label{section:region-routing-design-revised}

Định tuyến theo vùng nội dung xử lý vấn đề một trang PDF có thể chứa nhiều loại
nội dung cùng lúc. Hệ thống phát hiện các vùng văn bản, bảng, ảnh hoặc vùng
hình học trên trang, gán loại vùng và định tuyến sang bộ xử lý phù hợp. Đối với
văn bản, hệ thống dùng bộ trích xuất văn bản; đối với bảng, hệ thống dùng bộ
trích xuất bảng hoặc Hybrid TATR; đối với ảnh hoặc vùng scan, hệ thống dùng
OCR khi cần \cite{du2020ppocr}.

\begin{figure}[H]
\centering
\small
\begin{tabular}{c}
\fbox{\parbox{0.82\linewidth}{\centering Trang PDF}}\\[0.15cm]
\(\downarrow\)\\[-0.05cm]
\fbox{\parbox{0.82\linewidth}{\centering Phát hiện vùng: văn bản, bảng, ảnh, vùng hình học}}\\[0.15cm]
\(\downarrow\)\\[-0.05cm]
\fbox{\parbox{0.82\linewidth}{\centering Phân loại vùng: đoạn văn, tiêu đề, bảng, ảnh, chú thích}}\\[0.15cm]
\(\downarrow\)\\[-0.05cm]
\fbox{\parbox{0.82\linewidth}{\centering Tái dựng thứ tự đọc: bbox, bố cục một cột/hai cột}}\\[0.15cm]
\(\downarrow\)\\[-0.05cm]
\fbox{\parbox{0.82\linewidth}{\centering Định tuyến bộ xử lý: văn bản, bảng/Hybrid TATR, OCR}}\\[0.15cm]
\(\downarrow\)\\[-0.05cm]
\fbox{\parbox{0.82\linewidth}{\centering Block có cấu trúc: nội dung, trang, bbox, loại block, bộ xử lý}}\\
\end{tabular}
\caption{Luồng định tuyến theo vùng nội dung trong tầng ingest}
\label{fig:chapter4-region-routing-flow}
\end{figure}

Điểm quan trọng của thiết kế này là hệ thống xử lý từng vùng theo bản chất của
vùng đó, thay vì áp dụng một extractor cho toàn bộ trang. Điều này đặc biệt hữu
ích với PDF bán cấu trúc, nơi văn bản, bảng và hình ảnh thường xuất hiện xen kẽ.

\begin{table}[H]
\centering
\small
\renewcommand{\arraystretch}{1.2}
\caption{Lợi ích và đánh đổi của định tuyến theo vùng nội dung}
\label{tab:region-routing-tradeoff}
\begin{tabularx}{\textwidth}{
    |>{\raggedright\arraybackslash}p{0.24\textwidth}
    |>{\raggedright\arraybackslash}X
    |>{\raggedright\arraybackslash}X|
}
\hline
\textbf{Khía cạnh} & \textbf{Lợi ích} & \textbf{Đánh đổi / rủi ro} \\
\hline
Nội dung hỗn hợp &
Văn bản, bảng và ảnh được đưa sang bộ xử lý phù hợp thay vì xử lý toàn trang
bằng một phương pháp duy nhất. &
Nếu phát hiện hoặc phân loại vùng sai, bộ xử lý phía sau cũng có thể xử lý sai. \\
\hline
Truy vết nguồn &
Mỗi block giữ được trang, bbox, loại vùng và bộ xử lý đã dùng, giúp tạo dẫn chứng
chi tiết hơn. &
Metadata phức tạp hơn và cần chuẩn hóa thống nhất giữa các bộ xử lý. \\
\hline
Xử lý bảng &
Bảng được tách riêng để tái dựng hàng, cột và ô, hỗ trợ table-aware chunking. &
Bảng không có đường kẻ rõ, có ô gộp hoặc nhiều cấp tiêu đề vẫn có thể lỗi. \\
\hline
Hiệu năng &
Có thể tránh chạy OCR hoặc mô hình bảng trên toàn bộ trang khi không cần thiết. &
Bước phát hiện vùng và định tuyến tạo thêm chi phí xử lý ban đầu. \\
\hline
Fallback &
Khi một bộ xử lý không tạo được block hợp lệ, hệ thống có thể chuyển sang bộ xử
lý khác. &
Cần quy tắc kiểm tra chất lượng đầu ra để quyết định khi nào fallback. \\
\hline
\end{tabularx}
\end{table}

\section{Tái dựng bảng trong luồng ingest}
\label{section:chapter4-table-reconstruction}

Bảng được xử lý như một trường hợp riêng trong luồng ingest. Hệ thống không
chỉ lấy văn bản trong vùng bảng, mà cố gắng tái dựng quan hệ hàng, cột và ô. Đầu
ra của bước này là table block có bảng Markdown, danh sách cell và metadata
phục vụ table-aware chunking.

\begin{figure}[H]
\centering
\small
\begin{tabular}{c}
\fbox{\parbox{0.82\linewidth}{\centering Vùng bảng từ bước định tuyến vùng}}\\[0.15cm]
\(\downarrow\)\\[-0.05cm]
\fbox{\parbox{0.82\linewidth}{\centering Backend bảng: Default extractor, TATR-only hoặc Hybrid TATR}}\\[0.15cm]
\(\downarrow\)\\[-0.05cm]
\fbox{\parbox{0.82\linewidth}{\centering Nhận diện cấu trúc: hàng, cột, ô, bbox}}\\[0.15cm]
\(\downarrow\)\\[-0.05cm]
\fbox{\parbox{0.82\linewidth}{\centering Gán văn bản: hộp từ từ lớp văn bản PDF vào các ô; OCR là nguồn bổ sung khi tích hợp riêng}}\\[0.15cm]
\(\downarrow\)\\[-0.05cm]
\fbox{\parbox{0.82\linewidth}{\centering Block bảng: Markdown, danh sách ô, mã bảng, trang, bbox, bộ xử lý}}\\
\end{tabular}
\caption{Luồng tái dựng bảng trong tầng ingest}
\label{fig:chapter4-table-flow}
\end{figure}

Default extractor chủ yếu dựa trên text layer và quy tắc hình học nên nhanh, nhưng
có thể sai với bảng phức tạp. TATR-only dùng tín hiệu thị giác để nhận diện cấu
trúc bảng \cite{smock2022pubtables}, nhưng không tự gán đầy đủ nội dung chữ
vào từng ô. Trong triển khai hiện tại, Hybrid TATR kết hợp cấu trúc hình học từ
TATR với các hộp từ lấy trực tiếp từ lớp văn bản PDF. Nếu vùng bảng không có
hộp từ và chế độ chỉ dùng hình học không được bật, hệ thống giữ kết quả của bộ
trích xuất bảng ổn định làm phương án dự phòng. OCR được xem là nguồn hộp từ
bổ sung trong hướng tích hợp riêng cho tài liệu scan, không được gọi trực tiếp
bên trong module Hybrid TATR. Chi tiết định lượng của các bộ xử lý được trình
bày ở chương thực nghiệm.

\section{Chunking và indexing sau ingest}
\label{section:chunking-indexing-after-ingest}

Sau khi có block, hệ thống chuyển dữ liệu thành chunk để phục vụ truy xuất. Với
văn bản thông thường, chunk được tạo theo thứ tự đọc và ngữ cảnh đề mục. Với
bảng, hệ thống sinh thêm các dạng chunk có nhận thức bảng như table summary,
table structure, table row và table cell.

\begin{table}[H]
\centering
\small
\renewcommand{\arraystretch}{1.2}
\caption{Các dạng chunk sinh ra sau ingest}
\label{tab:ingest-chunk-types}
\begin{tabularx}{\textwidth}{
    |>{\raggedright\arraybackslash}p{0.24\textwidth}
    |>{\raggedright\arraybackslash}X
    |>{\raggedright\arraybackslash}p{0.26\textwidth}|
}
\hline
\textbf{Loại chunk} & \textbf{Dữ liệu chứa trong chunk} & \textbf{Vai trò} \\
\hline
Text chunk &
Nội dung đoạn văn, trang, heading path, block type. &
Truy xuất đoạn văn bản thông thường. \\
\hline
Table summary &
Caption, trang, số hàng/cột và tiêu đề chính. &
Tìm đúng bảng liên quan. \\
\hline
Table structure &
Biểu diễn bảng dạng Markdown hoặc cấu trúc tương đương. &
Giữ quan hệ tổng thể của bảng. \\
\hline
Table row &
Một hàng bảng cùng các cặp cột--giá trị. &
Trả lời câu hỏi theo một đối tượng hoặc một điều kiện. \\
\hline
Table cell &
Tiêu đề hàng, tiêu đề cột, giá trị ô và citation target. &
Truy xuất và dẫn chứng đến đúng ô. \\
\hline
\end{tabularx}
\end{table}

Indexing là bước xây dựng cấu trúc tìm kiếm từ các chunk. Với BM25, hệ thống
lưu thông tin phục vụ so khớp từ khóa~\cite{robertson2009bm25}. Với dense
retrieval, hệ thống lưu embedding của chunk để tìm kiếm theo độ tương đồng ngữ
nghĩa~\cite{karpukhin2020dpr,reimers2019sentencebert}. Các metadata như trang,
loại block, bảng, hàng và ô được giữ kèm chunk để tầng truy vấn có thể lọc, ưu
tiên hoặc tạo citation.

\section{Tổ chức mã nguồn của luồng ingest}
\label{section:ingest-code-organization}

\begin{table}[H]
\centering
\small
\renewcommand{\arraystretch}{1.2}
\caption{Các thành phần mã nguồn chính của luồng ingest}
\label{tab:ingest-code-organization}
\begin{tabularx}{\textwidth}{
    |>{\raggedright\arraybackslash}p{0.36\textwidth}
    |>{\raggedright\arraybackslash}X|
}
\hline
\textbf{Thành phần} & \textbf{Vai trò} \\
\hline
\texttt{app/ingest/probe.py} &
Thăm dò đặc trưng PDF và gợi ý chế độ ingest. \\
\hline
\texttt{app/ingest/region} &
Phát hiện region, phân loại region, vẽ overlay debug và định tuyến backend. \\
\hline
\texttt{app/ingest/extract} &
Các backend trích xuất text, layout, OCR, bảng và region-routed extraction. \\
\hline
\texttt{app/ingest/reading\_order.py} &
Sắp xếp block theo thứ tự đọc dựa trên hình học trang. \\
\hline
\texttt{app/ingest/table\_chunking.py} &
Sinh table summary, table structure, table row và table cell chunk. \\
\hline
\texttt{app/ingest/pipeline.py} &
Điều phối ingest, validation, fallback và chuẩn hóa đầu ra. \\
\hline
\end{tabularx}
\end{table}

\section{Tổng kết chương}
\label{section:chapter4-ingest-summary}

Chương này đã trình bày luồng ingest và trích xuất PDF của hệ thống. Thiết kế
được tổ chức quanh ba ý chính: thăm dò tài liệu để chọn chiến lược ingest, định
tuyến theo vùng để xử lý các loại nội dung khác nhau, và giữ lại dữ liệu có cấu
trúc cùng metadata nguồn. Kết quả của chương này là tập block và chunk có
metadata, làm đầu vào cho bước xây dựng chỉ mục truy xuất và luồng xử lý truy
vấn ở chương tiếp theo.

\end{document}
